
\begin{proposition}[Eliminating vertical conductor on a smooth ruled cusp normalization]\label{cdp:st:A.5.10}
Let $Y$ be a smooth compact complex threefold with
$H^2(Y,\mathbb Q)=0$, let $S\subset Y$ be a reduced
irreducible divisor, and suppose its normalization
$h:Z\to S\subset Y$ is a smooth irrational
ruled surface $p:Z\to B$.
Suppose its conductor divisor is $D_Z=2C+V$, where
$C$ is a ruling section and $V$ is vertical, and
the generic transversal along $h(C)$ is an ordinary cusp.
Then $B$ is elliptic, $V=0$, $C^2=0$, and $Z$ is
already minimal.
The image $h(C)$ is smooth elliptic and $S$ is locally
an ordinary-cusp product along it. No local
complete-intersection hypothesis on the image curve is
required. The conductor-image scheme along the vertical
components may initially be nonreduced.
\end{proposition}
\begin{proof}
Let $g=g(B)\geq1$, let $F$ be the full ruling fibre
class, and write $V=\sum_i c_i\Gamma_i$ for its
irreducible vertical components. Put
$v=V\cdot C$, $n=C^2$ and $k=2g-2$.
Conductor duality and Cartier adjunction give the exact
line-bundle comparison
\[
 K_Z\otimes\mathcal O_Z(2C+V)
       =h^*(K_Y\otimes\mathcal O_Y(S)).
                                                        \tag{A.J.21}\label{cdp:eq:A.J.21}
\]
Its right side has rational first Chern class zero.
Intersecting with $C$ and using adjunction on the section
gives
\[
 n=2-2g-v=-k-v.                                  \tag{A.J.22}\label{cdp:eq:A.J.22}
\]
Let $D_p$ be the divisorial common zero of
$dp:T_Z\to p^*T_B$. Its remaining zero scheme is
zero-dimensional, say $\Xi$, so the image is
$I_\Xi\otimes p^*T_B(-D_p)$.
The kernel $T_p$ is an invertible sheaf, by the depth
lemma. Taking determinants gives the actual comparison
\[
 T_p=K_Z^{-1}\otimes p^*K_B\otimes\mathcal O_Z(D_p).
\]
The section makes $dp$ surjective at every point of $C$.
Consequently $D_p$ and $\Xi$ are disjoint from $C$,
and $T_p|_C=N_{C/Z}$.

Put $\mathcal F=h^*T_Y$. Poincare duality on the compact
oriented six-manifold $Y$ also gives
$H^4(Y,\mathbb Q)=0$. Therefore
\[
 c_1(\mathcal F)=0,\qquad c_2(\mathcal F)=0
        \quad\hbox{in rational cohomology}.       \tag{A.J.23}\label{cdp:eq:A.J.23}
\]
The differential is generically injective. Let
$L\subset\mathcal F$ be the saturation of the image
of $T_p\xrightarrow{dh}\mathcal F$.
It is invertible: a saturated rank-one subsheaf of a
locally free sheaf on a smooth surface is reflexive.
The map $T_p\to L$ has no vertical divisorial zero.
Indeed $h$ is finite, so no ruling fibre is contracted;
in characteristic zero its derivative along that fibre
is nonzero generically.

At the generic point of $C$, an ordinary cusp has
normalization differential equivalent to that of
$(u,s)\mapsto(u^2,u^3,s)$.
The ruling tangent either lies in its kernel along $C$
or does not. The vanishing order in the first case is
exactly one, and in the second case it is zero.
Write $\epsilon=1$ or $\epsilon=0$ for these respective
alternatives. There are no other horizontal conductor
components, and the normalization is an immersion away
from the conductor. Hence the actual line comparison is
\[
 L=T_p\otimes\mathcal O_Z(\epsilon C).
                                                        \tag{A.J.24}\label{cdp:eq:A.J.24}
\]
This checks both positions of the derivative kernel
relative to the ruling tangent.

The quotient $E=\mathcal F/L$ is torsion-free of rank
two. The differential induces a generically nonzero map
\[
 T_Z/T_p=I_\Xi\otimes p^*T_B(-D_p)
                    \longrightarrow E\longrightarrow E^{**}.
\]
It extends uniquely to a map from $p^*T_B(-D_p)$ into
$E^{**}$, since both are locally free and the omitted
set has codimension two; locally this is Hartogs
extension of its coefficient functions.
Let $M$ be its saturated image line in the rank-two
bundle $E^{**}$. At the generic point of $C$, the
two-by-two differential minors have common vanishing
order one. Removing the order $\epsilon$ already
assigned to $L$ leaves order $1-\epsilon$ for this map.
At the generic point of $\Gamma_i$, take a plane transversal
to its image and a normalization parameter $u$ for the
corresponding plane-curve branch. If $m_i$ is its
multiplicity, its derivative modulo the nonzero tangent
along $\Gamma_i$ vanishes to order $r_i=m_i-1$.
These are all codimension-one zeros. Put
$R=\sum_i r_i\Gamma_i$ and $r=R\cdot C$. We obtain
\[
 M=p^*T_B(-D_p)\otimes
           \mathcal O_Z\bigl((1-\epsilon)C+R\bigr).
                                                        \tag{A.J.25}\label{cdp:eq:A.J.25}
\]

The plane-curve conductor gives the inequalities
\[
 2R\leq V,\qquad 2r\leq v.                      \tag{A.J.26}\label{cdp:eq:A.J.26}
\]
To spell out their local length content, on a single
plane branch of multiplicity $m_i$ the classes
$u,u^2,\ldots,u^{m_i-1}$ are linearly independent in
the normalization quotient. Its delta invariant is thus
at least $m_i-1$. The branch is Gorenstein. Finite
duality identifies the length of its ring modulo its
conductor with the length of its normalization quotient;
the intrinsic conductor exponent is consequently twice
that delta invariant. Other branches add their
nonnegative intersection multiplicities to this exponent.
For the last comparison, writing the plane equation
as $\prod_j f_j$ multiplies the residue denominator on
branch $i$ by $\prod_{j\ne i}f_j$, whose valuation is
exactly the sum of those intersection multiplicities.
This proves \eqref{cdp:eq:A.J.26} for every branch, including cases
with nonreduced conductor-image scheme.

Let $\ell_1$ be the length of $E^{**}/E$. Saturation
of $M$ gives a zero-dimensional subscheme $\Xi_M$ and a
line bundle $N$ with
\[
 0\longrightarrow M\longrightarrow E^{**}
       \longrightarrow I_{\Xi_M}\otimes N\longrightarrow0.
\]
Write $\ell_2=\operatorname{length}\Xi_M$.
First Chern classes give $c_1(N)=-c_1(L)-c_1(M)$.
The second Chern classes in these typed quotient
sequences and \eqref{cdp:eq:A.J.23} give the intersection identity
\[
 L^2+L\cdot M+M^2=\ell_1+\ell_2\geq0.            \tag{A.J.27}\label{cdp:eq:A.J.27}
\]
Indeed $c_2(E)=L^2$,
$c_2(E^{**})=L^2-\ell_1$, and
$c_2(E^{**})=-L\cdot M-M^2+\ell_2$.

Set $W=V+D_p$ and $B_v=R-D_p$, both vertical
divisors, the first effective. The exact comparisons
above and \eqref{cdp:eq:A.J.21} give the numerical classes
\[
 L\equiv(2+\epsilon)C+W+kF,\qquad
 M\equiv(1-\epsilon)C+B_v-kF.
\]
The intersection form on vertical divisors is negative
semidefinite. One can check this directly by contracting
the ruling blowups: every vertical class is a multiple
of the pullback fibre class plus exceptional classes,
on which the form is negative definite. Hence
\[
 q_v=W^2+W\cdot B_v+B_v^2
      =(W+B_v/2)^2+3B_v^2/4\leq0.
\]
Its vanishing implies $W^2=B_v^2=0$; these classes
are then numerical multiples of the fibre class.

For $\epsilon=1$, using $D_p\cdot C=0$ in the
left side of \eqref{cdp:eq:A.J.27} gives
\[
 \ell_1+\ell_2=-6k-3v+3r+q_v
                  \leq-6k-3v/2.
\]
Nonnegativity, with $k\geq0$ and $v\geq0$,
forces $g=1$, $v=0$, $r=0$ and $q_v=0$.

For $\epsilon=0$, the same intersection calculation reads
\[
 0\leq\ell_1+\ell_2
       =-6k-2v+4r+q_v\leq-6k.                   \tag{A.J.28}\label{cdp:eq:A.J.28}
\]
It forces $g=1$, $\ell_1=\ell_2=0$, $2r=v$,
and $q_v=0$.
Thus $E=E^{**}$ is locally free, and
$L\to\mathcal F$ is a subbundle inclusion.
Along $C$, the composite
$T_C\to\mathcal F|_C\to E|_C$ vanishes generically,
because $dh$ has rank one there and its nonzero ruling
tangent image is $L|_C$. It therefore vanishes
everywhere. The exact subbundle sequence gives an
actual nonzero holomorphic line-bundle map
\[
 T_C\longrightarrow L|_C=T_p|_C=N_{C/Z}.
                                                        \tag{A.J.29}\label{cdp:eq:A.J.29}
\]
It is nonzero since the finite map $C\to h(C)$ has
nonzero derivative generically. Its Hom line has degree
$n-(2-2g)=-v$ by \eqref{cdp:eq:A.J.22}. A nonzero section on a
compact curve has nonnegative degree, so $v=0$ in this
alternative too.

In both alternatives, $v=r=0$ and $q_v=0$.
Thus $W$ is numerically a fibre multiple and
$W\cdot C=V\cdot C+D_p\cdot C=0$, so it is
numerically zero. An effective nonzero divisor on the
projective ruled surface has positive intersection
with an ample class. It follows that $W=0$ and hence
$V=D_p=0$. Now $B_v=R$ is effective, has square zero
and has zero intersection with $C$, so the same
argument gives $R=0$. Equation \eqref{cdp:eq:A.J.22} gives $C^2=0$.
Finally $K_Z+2C$ is numerically trivial. A ruling
$(-1)$-curve $E_0$ would satisfy
$0=(K_Z+2C)\cdot E_0=-1+2C\cdot E_0$, impossible
because $C\cdot E_0$ is an integer. There are therefore
no ruling blowups and $Z$ is minimal.

Statement \ref{cdp:st:A.5.7} now applies everywhere along
the conductor: the normalization is smooth and its
conductor is exactly $2C$. It makes the image smooth
and the surface locally the cusp product. The image is
elliptic because $C\to h(C)$ is an isomorphism.
\end{proof}


